\documentclass[border=10pt]{standalone}
\usepackage{amsmath,amssymb}
\usepackage{xcolor}
\usepackage{tcolorbox}
\tcbuselibrary{skins}

\pagecolor{yellow!10}
\definecolor{questioncolor}{RGB}{0,100,0}

\begin{document}

\begin{tcolorbox}[
  enhanced,
  colback=white,
  colframe=questioncolor,
  boxrule=1pt,
  arc=5pt,
  outer arc=5pt,
  drop shadow={opacity=0.5},
  fontupper=\color{questioncolor}\sffamily
]
During a hostage situation in a 10m $\times$ 10m $\times$ 3m room, tear gas is released from a canister at the center of the floor. The concentration of tear gas $c(x,y,t)$ can be modeled by the 2D heat equation:

\begin{equation*}
\frac{\partial c}{\partial t} = D\left(\frac{\partial^2 c}{\partial x^2} + \frac{\partial^2 c}{\partial y^2}\right)
\end{equation*}

where $D = 0.05$ m²/s is the diffusion coefficient. If the initial concentration at the point of release $(5,5)$ is 100 mg/m³ and zero elsewhere, what is the concentration at point $(2,2)$ after 10 seconds?
\end{tcolorbox}

\end{document}